\documentclass{ieeeaccess}
\usepackage{graphicx,amsmath,amssymb,bm,booktabs,array,textcomp,url}
\usepackage{adjustbox}
\usepackage[hidelinks]{hyperref}
\usepackage[numbers,sort&compress]{natbib}
\usepackage{microtype}
\usepackage[htt]{hyphenat}
\sloppy
\emergencystretch=3em
\providecommand{\tightlist}{\setlength{\itemsep}{0pt}\setlength{\parskip}{0pt}}
\providecommand{\passthrough}[1]{#1}
\providecommand{\pandocbounded}[1]{#1}
\providecommand{\real}[1]{#1}

\history{Date of submission June 8, 2026.}
\doi{10.1109/ACCESS.2026.DOI}

\title{Cache the Known, Ask Only the Novel: A Confidence-Routed Harness for Sub-1B Language Models}
\author{\uppercase{Ryota Nakamura}\authorrefmark{1}}
\address[1]{Faculty of Data Science, Musashino University, Tokyo 135-8181, Japan (e-mail: ryota.nakamura@ds.musashino-u.ac.jp)}
\corresp{Corresponding author: Ryota Nakamura (e-mail: ryota.nakamura@ds.musashino-u.ac.jp).}
\markboth{R. Nakamura: Cache the Known, Ask Only the Novel}{R. Nakamura: Cache the Known, Ask Only the Novel}

\begin{document}
\maketitle

\begin{abstract}
On-device deployment increasingly favors sub-1-billion-parameter ("tiny") language models for their low latency, privacy, and cost, yet such frozen models are unreliable on structured classification. The usual response is to make the model better (fine-tuning, larger models). We argue the opposite for tiny-model systems: do not strengthen the LLM, but **separate the known from the novel** and route each to where it is cheapest and most reliable. We instantiate this as a deterministic harness around a **frozen** sub-1B model with an **experience cache** of verified (text, label) cases, exposing two complementary paths --- a non-parametric predictor (embedding kNN / logistic regression) that never calls the LLM, and the LLM conditioned on retrieved exemplars and distilled rules --- plus a **confidence gate** that decides between them. The central engineering result is that this gate is principled, not heuristic: the non-parametric classifier's own confidence predicts whether its (no-LLM) answer is correct at mean AUROC 0.83 across seven datasets, whereas nearest-neighbor similarity alone is near chance. Routing by that confidence matches or beats the better single path (85.6% vs 83.9% no-LLM, 57.0% LLM) while sending only ~4% of inputs to the LLM, and exceeds both where they are complementary (e.g., 79/91$\rightarrow$97% on Japanese sentiment), approaching an oracle ceiling of 92.1%. A comprehensive sweep (11 datasets $\times$ 8 sub-1B models, 88 cells) shows the no-LLM path dominates (71.9% mean vs 26--29% for every LLM-path arm), is nearly invariant to class count ($\approx$80% from 2 to 77 classes), and matches LoRA at zero training cost; structuring the model's own failures into discriminative rules helps the LLM path but only marginally on average, and beats structuring its successes. The LLM remains essential for genuine novelty: on unseen-theme (leave-one-theme-out) queries it generalizes (92%) where the non-parametric path drops (56%). Findings are validated with leave-one-theme-out splits, label-corruption and 1-NN controls, paired McNemar tests, Wilson confidence intervals, and an external re-test on six real-world datasets. The contribution is thus a *system-centric* design principle for sub-1B **classification** deployments --- delegate known inputs to deterministic components and spend the LLM only on the genuinely novel --- with a validated routing rule that makes it operational. We deliberately scope all claims to closed-set classification, whose outputs are recoverable from verified cases; a generative-QA pilot marks the boundary where the no-LLM path stops transferring, and extending the principle to open-ended generation and reasoning is left to future work.
\end{abstract}

\begin{IEEEkeywords}
Small language models, on-device inference, experience cache, retrieval-augmented classification, query routing, selective prediction, cost-efficient inference, edge AI.
\end{IEEEkeywords}

\input{_body_ieee.tex}

\bibliographystyle{IEEEtran}
\bibliography{references}
\end{document}
